\slajdid{09_3-s006}
\begin{frame}[label=09_3-s006]
\gusttekst\setlength{\abovedisplayskip}{3pt}\setlength{\belowdisplayskip}{3pt}\setlength{\abovedisplayshortskip}{2pt}\setlength{\belowdisplayshortskip}{2pt}\textbullet\quad\alert{Optimizacija kašnjenja u jednom CMOS invertoru}
\par\medskip
Analizirajući statičke karakteristike CMOS invertora već smo videli da bi zbog različitih pokretljivosti nosilaca i dimenzije P i N tranzistora trebale da budu različite, da bi dobili „dobru“ karakteristiku prenosa. Na žalost to će uticati na ukupnu kapacitivnost kojom je invertor opterećen ako mu se na izlazu nalazi isti takav invertor, odnosno na kapacitivnost kojom invertor opterećuje prethodno stepen. Zbog toga ćemo posmatrati realnu situaciju koja nam se u digitalnom sistemu pojavljuje, veze dva invertora
\par\medskip\centering\begin{tikzpicture}[x=.8cm,y=.65cm,font=\sitneoznake,>=Latex]
\foreach \x/\i in{0/1,3.8/2}{\begin{scope}[shift={(\x,0)}]
\draw(0,2.65)--(0,3.35);\draw(-.13,2.6)--(-.13,3.4);\draw(-.23,3)circle(.1);\draw(-.33,3)--(-1,3)--(-1,.8)--(-.13,.8);
\draw(0,3.3)--(.4,3.3)--(.4,3.9);\draw(.15,3.9)--(.65,3.9)node[above left]{$V_{DD}$};\draw(0,2.7)--(.4,2.7)--(.4,1.1);
\draw(0,.45)--(0,1.15);\draw(-.13,.4)--(-.13,1.2);\draw(0,1.1)--(.4,1.1);\draw(0,.5)--(.4,.5)--(.4,-.1)node[ground]{};
\node[right]at(.55,3){P\i};\node[right]at(.55,.8){N\i};\end{scope}}
\draw(-1.65,1.9)node[ocirc]{}node[left]{$V_{I1}$}--(-1,1.9);\draw(.4,1.9)--(2.8,1.9);\node[above right]at(.4,1.9){$V_{O1}$};\node[above left]at(2.8,1.9){$V_{I2}$};\draw(4.2,1.9)--(5.1,1.9)node[ocirc]{}node[right]{$V_{O2}$};
\end{tikzpicture}

\end{frame}
